\chapter{设备驱动、中断下半部与 DMA}

\section{异步设备请求与完成证据}

设备驱动把硬件寄存器、中断、DMA 和内核对象连接起来。应用看到的是 fd、socket、块设备或帧缓冲；驱动看到的是 MMIO 寄存器、描述符环、设备状态位、中断线、DMA 地址和错误码。驱动的难点在于硬件和 CPU 并发运行：CPU 写描述符，设备异步读取；设备写内存，CPU 稍后读取；中断可能在任意时刻到来。

先从一个错误直觉开始：

\begin{lstlisting}
disk_write(buf, len):
  return disk.write(buf, len)
\end{lstlisting}

这像普通函数调用，但真实设备不会沿着你的调用栈返回。CPU 能做的是写设备寄存器、填一段内存中的描述符、敲 doorbell 通知设备；设备随后按自己的时钟和状态机工作。它可能成功，可能失败，可能很久以后才完成，也可能在 reset 中丢掉请求。驱动的工作就是把这种异步硬件行为包装成内核可管理的请求。

把最小路径写清楚：

\begin{lstlisting}
submit_write(buf):
  req = allocate_request()
  req.buffer = buf
  req.state = PREPARED
  dma = map_buffer_for_device(buf)
  fill_descriptor(req.id, dma, len)
  publish_descriptor_with_barrier()
  ring_doorbell()
  req.state = IN_FLIGHT
  sleep_until_completion(req)
\end{lstlisting}

这段伪代码第一次引出几个概念。descriptor 是给设备看的请求摘要，通常包含 DMA 地址、长度、标志和请求身份。doorbell 是通知设备“我放了新活”的寄存器写入。completion 是设备稍后交回的完成结果。request 是内核自己的对象，用来把用户请求、buffer、描述符、等待任务和最终错误码连在一起。

最简单设备可以用 programmed I/O：CPU 直接读写设备寄存器或数据端口。PIO 易懂，但吞吐低，CPU 被数据搬运占住。DMA 让设备直接读写内存，CPU 只准备描述符和缓冲区。DMA 提高吞吐，也引入缓存一致性、地址映射、内存屏障和 buffer 生命周期问题。

PIO 和 DMA 的区别不是“老技术”和“新技术”的区别，而是 CPU 参与程度不同：

\begin{longtable}{p{0.18\linewidth}p{0.34\linewidth}p{0.34\linewidth}}
\toprule
方式 & CPU 做什么 & 风险边界 \\
\midrule
PIO & 每次读写设备寄存器搬运数据 & CPU 忙等、吞吐低、长时间占用核心 \\
DMA & 准备 buffer 和描述符，让设备直接搬运内存 & buffer 生命周期、cache 一致性、IOMMU 映射 \\
中断 & 设备完成后通知 CPU & 中断风暴、上半部不能睡眠、completion 要能找回请求 \\
轮询 & CPU 定期检查完成队列 & 延迟/吞吐折中、budget 背压 \\
\bottomrule
\end{longtable}

驱动通常分成上半部和下半部。中断上半部快速确认设备、读取状态、屏蔽或清除中断，把复杂工作推迟到底半部、工作队列、软中断或普通内核线程。这样可以减少中断关闭时间，避免在硬中断上下文执行可能睡眠的操作。

驱动和前面章节的连接也要明确：提交请求时可能让当前 task 睡眠，completion 到来后要唤醒 wait queue；DMA buffer 来自页分配器和虚拟内存系统；设备错误最终要通过 fd、socket、块层或文件系统返回给用户；权限和 namespace 决定用户能否打开这个设备。驱动不是孤立章节，而是硬件异步事件进入 OS 对象系统的入口。

\section{从发现设备到安全移除}

\subsection{为什么设备和驱动要成为两个对象}

若内核为每一块具体硬件直接写一段固定初始化代码，单一机器可以启动，但设备替换、热插拔、
休眠恢复和同类多设备都会使这套做法失效。内核必须分别保存两类事实：机器上实际出现了
哪些设备，以及当前有哪些软件能够控制某类设备。前者形成 device 对象，后者形成 driver 对象。

两者通过可匹配的身份信息建立绑定。例如总线发现一个设备后，先登记设备身份、资源范围和
父子关系；驱动注册时声明自己支持哪些身份。匹配成功只说明“可以尝试控制”，还不能立即向
用户发布设备，因为初始化过程中仍可能缺少中断、DMA 队列或电源资源。

\begin{lstlisting}
discover device
  -> create device object
  -> match a compatible driver
  -> probe and acquire resources
  -> initialize queues and interrupt handling
  -> publish usable interface
\end{lstlisting}

这个顺序使失败保持局部。若初始化队列失败，驱动释放此前取得的资源，device 对象仍可保留
为“已发现但尚不可用”，也可以等待另一个驱动；如果在资源尚未齐备时就提前创建设备节点，
用户请求可能进入一个半初始化对象。

\subsection{probe 为什么是一项可回滚事务}

一次初始化通常按顺序取得寄存器映射、中断向量、DMA 能力、队列内存、设备固件和上层注册。
任一步都可能失败。资源释放必须采用相反顺序，因为较后的对象往往引用较早的对象：队列使用
DMA 映射，完成路径使用中断，用户接口引用队列和设备状态。

\begin{lstlisting}
probe(dev):
  regs = map_registers(dev)
  irq = allocate_interrupt(dev)
  queues = allocate_dma_queues(dev)
  reset_and_negotiate(dev)
  start_device(dev, queues)
  publish_interface(dev)
  return success

on failure:
  unpublish_if_needed()
  stop_device_if_started()
  free_dma_queues()
  release_interrupt()
  unmap_registers()
\end{lstlisting}

这里的提交点是向其他子系统发布可用接口。在此之前，错误只需释放私有资源；在此之后，
关闭和移除还必须阻止新请求，并处理已经取得对象引用的调用者。托管资源可以减少遗漏，
但不会替驱动决定设备是否仍在 DMA、请求是否已完成以及用户接口何时不可见。

\subsection{移除为何不能直接释放对象}

设备被拔出或驱动卸载时，仍可能存在三类活动：尚未提交到设备的新请求、设备已经接收的
in-flight 请求，以及持有 fd 或其他引用的用户。若移除路径直接释放队列和 DMA buffer，
迟到的中断或设备写入就可能访问已经复用的内存。

安全移除通常先把设备状态改为 quiescing，使新请求立即失败或留在上层；随后停止队列、
屏蔽并同步中断，等待设备不再访问 DMA 映射；未完成请求以明确错误完成；最后撤销用户可见
接口并等待对象引用归零。只有这些条件同时成立，寄存器、队列和设备对象才能释放。

\begin{longtable}{@{}F{0.19\linewidth}F{0.38\linewidth}F{0.33\linewidth}@{}}
\toprule
\rowcolor{BookBlueTint}
阶段 & 必须建立的事实 & 过早进入下一阶段的后果 \\
\midrule
停止接收 & 新请求不能再进入设备队列 & 关闭过程永远追不上新提交 \\\tablerowrule
终止设备访问 & 设备不再读取描述符或 DMA buffer & 设备访问已释放内存 \\\tablerowrule
同步完成路径 & 旧中断和延后工作已经退出 & 迟到完成使用旧 request 或 tag \\\tablerowrule
完成未决请求 & 所有等待者收到成功或错误 & 线程永久睡眠，引用无法归零 \\\tablerowrule
撤销与释放 & 接口不可见且外部引用为零 & 用户通过旧 fd 进入已销毁对象 \\
\bottomrule
\end{longtable}

复位是同一问题的较小版本。复位前后的请求 tag 可能数值相同，因此驱动还需要 generation
或等价状态区分旧队列和新队列。完成记录只有同时匹配队列世代和活动请求，才能成为释放
buffer 的证据。

\section{队列所有权与中断分工}

先画一个网卡或磁盘的描述符环：

\begin{lstlisting}
descriptor ring:
  head: device consumes descriptors
  tail: driver produces descriptors
  desc[i]:
    dma_address
    length
    flags
    status
\end{lstlisting}

驱动发送路径通常是：从内核分配 buffer，映射成设备可见 DMA 地址，填描述符，执行内存屏障，更新 tail，通知设备。完成路径通常是：设备 DMA 完成并写 status，触发中断，驱动回收描述符，解除 DMA 映射，释放 buffer，唤醒等待者或上报 completion。

\section{描述符环、DMA 与内存顺序}

\subsection{MMIO 寄存器访问}

MMIO 把设备寄存器映射到物理地址。访问 MMIO 不能当成普通内存优化：编译器不能随便重排，CPU 也要遵守设备要求的顺序。驱动通常使用专用 readl/writel 一类接口。

\begin{lstlisting}
device_start(dev):
  writel(dev->mmio + CTRL, CTRL_RESET)
  wait_until(readl(dev->mmio + STATUS) & READY)
  writel(dev->mmio + IRQ_MASK, RX_IRQ | TX_IRQ)
\end{lstlisting}

轮询 wait 要有超时。设备可能坏掉，寄存器可能永远不变。驱动若在内核里无限等待，会拖死系统。

\subsection{描述符环}

描述符环是设备驱动的核心队列。生产者和消费者可以是 CPU 和设备，也可以反过来。环满时不能继续提交，环空时不能继续回收。

\begin{lstlisting}
submit_tx(desc_ring, packet):
  next = (tail + 1) % ring_size
  if next == head:
    return -EAGAIN
  desc[tail].addr = dma_map(packet.data)
  desc[tail].len = packet.len
  desc[tail].flags = OWNED_BY_DEVICE
  memory_barrier()
  tail = next
  writel(DOORBELL, tail)
\end{lstlisting}

这里的屏障非常关键。CPU 必须先把描述符内容写完，再更新 tail 或 doorbell 通知设备；否则设备可能看到半初始化描述符。

把一次网卡发送请求走细一点。用户态 \code{send} 最终让内核拿到一段要发送的数据，网络栈构造 packet buffer，驱动要把这个 buffer 交给网卡。网卡不认识 \code{struct sk\_buff}，也不能读用户虚拟地址；它通常只读一段描述符内存，描述符里写着 DMA 地址、长度、校验和 offload 标志、结束位和 request tag。CPU 是生产者，设备是消费者。

\begin{longtable}{p{0.16\linewidth}p{0.42\linewidth}p{0.3\linewidth}}
\toprule
步骤 & 驱动和硬件动作 & 不能错的边界 \\
\midrule
保留槽位 & 驱动读取 \code{head/tail}，确认环未满，选择 \code{tail} 槽位 & 满环必须返回或停止上层队列，不能覆盖设备未消费的描述符 \\
映射 buffer & DMA API 把 packet buffer 映射成设备可见 DMA 地址 & 设备不能拿用户指针或内核虚拟地址 \\
填写描述符 & 写入地址、长度、flags、tag，最后设置 owned/valid 位 & valid 位不能早于其他字段对设备可见 \\
发布顺序 & 执行 DMA 写屏障，保证描述符内容先于 doorbell & 少屏障会让设备读到旧字段或半初始化字段 \\
敲门铃 & 写 MMIO doorbell 或更新 tail，通知设备取新描述符 & doorbell 只是通知，不代表发送完成 \\
设备 DMA & 网卡读取描述符和数据，经 IOMMU 访问内存，发送到链路 & buffer 在 completion 前不能释放或复用 \\
完成回收 & 设备写回 completion 或更新 head，驱动 unmap 并释放 buffer & 回收必须以设备完成证据为准 \\
\bottomrule
\end{longtable}

这里有两个常见误解。第一，\code{tail = next} 只是软件变量变化，设备未必看见；真正通知设备的通常是 MMIO doorbell 或共享队列尾指针的可见更新。第二，completion 不等于用户业务语义完成。网卡发送 completion 通常说明设备已经不再需要这个 DMA buffer，未必说明对端应用已经收到数据；块设备 completion 说明设备处理了请求，是否持久还要看 flush/FUA 和文件系统语义。驱动层必须先把“buffer 能否释放”这件事讲清楚，再谈上层协议。

迟到 completion 是取消路径最容易出错的地方。假设上层超时，驱动 reset 队列并把 request 对象释放；如果设备或旧中断稍后又报告同一个 tag 完成，而 tag 已经被新请求复用，驱动可能唤醒错任务或释放错 buffer。可靠实现通常给 request tag、队列 generation、in-flight bitmap 和 reset 状态建立明确规则：reset 开始后停止新提交，标记旧请求失败，等待或屏蔽旧 completion，确认设备不再 DMA 后再复用 tag 和 buffer。

\subsection{DMA 和缓存一致性}

若 CPU cache 和设备 DMA 不自动一致，驱动必须在设备读内存前清理 cache，在设备写内存后失效 cache。即使硬件支持 IOMMU，也要建立设备可访问地址映射。

\begin{lstlisting}
prepare_dma_to_device(buf):
  cache_clean(buf)
  dma_addr = iommu_map(buf.phys)
  return dma_addr

finish_dma_from_device(buf):
  cache_invalidate(buf)
  parse_received_data(buf)
\end{lstlisting}

DMA buffer 的生命周期必须覆盖设备访问全过程。提交后不能释放或复用 buffer，直到 completion 表明设备不再访问它。这个规则和异步 I/O 的 buffer 所有权完全一致。

缓存一致性要分方向看。设备从内存读数据，叫 TO\_DEVICE：CPU 先把数据写进 buffer，但这些写入可能还停在 cache 中，设备从内存读会看到旧内容；所以提交前要 clean 或通过 coherent 映射保证设备可见。设备向内存写数据，叫 FROM\_DEVICE：设备已经把新数据写进内存，但 CPU cache 里可能还有旧 cache line；所以 completion 后 CPU 读取前要 invalidate 或同步。

\begin{longtable}{p{0.18\linewidth}p{0.36\linewidth}p{0.34\linewidth}}
\toprule
方向 & 提交前/完成后要做什么 & 忘记后的现象 \\
\midrule
TO\_DEVICE & CPU 写 buffer 后，clean cache，再让设备 DMA 读 & 设备发送旧包、旧磁盘数据或半初始化描述符 \\
FROM\_DEVICE & 设备 completion 后，invalidate cache，再让 CPU 解析 & CPU 读到旧包内容，协议栈校验失败或数据错乱 \\
BIDIRECTIONAL & 两个方向都要同步，并严格定义所有权切换 & CPU 和设备同时修改，结果半新半旧 \\
coherent DMA & 平台保证 CPU/设备可见性，但仍需要顺序屏障 & 数据可见不等于 doorbell 顺序自动正确 \\
\bottomrule
\end{longtable}

再看一个接收包场景。驱动预先分配 RX buffer，把它 DMA map 成 FROM\_DEVICE，挂到 RX ring。挂上去之后，buffer 暂时归设备所有，CPU 不应把它当普通内存写。网卡收到包后 DMA 写入 buffer，并在描述符里写入长度、校验状态和完成位。驱动 poll 到完成位后，先执行 DMA 同步或 invalidate，再读取包头和长度，构造 skb 交给协议栈。若顺序反了，CPU 可能先读旧 cache line，看到旧长度或旧以太网头；这类 bug 往往不是每次复现，因为它取决于 cache 状态、CPU 迁移和包到达时机。

IOMMU 还提供了另一个边界：设备只能访问被映射的 IOVA 范围。驱动 unmap 后，设备若继续 DMA，正确结果应是 IOMMU fault 或设备错误，而不是静默覆盖任意物理内存。可这要求驱动先确认设备不再持有该 descriptor，再 unmap 并释放页。IOMMU 是安全网，不是生命周期管理的替代品。

\subsection{中断合并和轮询}

高吞吐设备如果每个包都中断一次，CPU 会被中断淹没。中断合并让设备积累多个完成再中断；NAPI 一类机制在高负载下从中断转为轮询，批量处理包，降低中断成本。代价是单个包延迟可能增加。

\begin{lstlisting}
rx_interrupt():
  disable_rx_irq()
  schedule_poll()

poll(budget):
  processed = 0
  while processed < budget and rx_desc_ready():
    handle_packet()
    processed += 1
  if no_more_packets():
    enable_rx_irq()
\end{lstlisting}

这里的 budget 是背压工具。一次 poll 不能无限处理，否则其他任务得不到 CPU。

\section{超时、复位与迟到完成}

\begin{itemize}
\tightlist
\item MMIO 等待必须有超时，设备无响应时返回错误。
\item 描述符环满时返回背压，不覆盖未完成描述符。
\item 提交描述符前执行必要内存屏障。
\item DMA buffer 在 completion 前不能释放或复用。
\item 中断处理上半部不执行可能睡眠操作。
\item 批量 poll 遵守 budget，不能垄断 CPU。
\item 设备 reset 后，驱动能清理 in-flight 请求并通知上层失败。
\end{itemize}

\subsection{设备完成如何支撑上层提交}

块层原本可以单独讲很久，但放在设备章里看更清楚：它是把文件系统的逻辑读写转换成设备队列 request 的中间层。文件系统提交的是 bio，描述“文件系统需要这些扇区读写”；blk-mq 把 bio 合并、排序、分派给硬件队列；驱动把 request 变成设备描述符；completion 再把错误和字节数一路传回等待者。

\begin{longtable}{p{0.18\linewidth}p{0.40\linewidth}p{0.30\linewidth}}
\toprule
层次 & 保存什么账本 & 不能混淆的边界 \\
\midrule
bio & 一组页、扇区范围、读写方向和完成回调 & bio 完成不一定等于文件系统事务提交 \\
request & 合并后的设备队列项、tag、超时和状态 & request tag 复用必须等旧 completion 收敛 \\
hardware queue & CPU/NUMA 亲和的提交队列 & 多队列提高吞吐，也引入负载和顺序问题 \\
flush/FUA & 写缓存顺序和持久化约束 & 普通写完成不等于掉电安全 \\
I/O scheduler & 合并、排序、公平和延迟控制 & 吞吐优化不能破坏屏障语义 \\
\bottomrule
\end{longtable}

存储顺序比普通异步完成更严格。日志文件系统、数据库和键值存储经常要求“先写数据，再提交元数据或 manifest”。如果设备或调度器为了吞吐随意重排，就可能在崩溃后看到新 manifest 指向旧数据。块层用 flush、FUA、barrier 或等价机制表达这些顺序约束；文件系统还要把自己的 journal commit、checkpoint 和目录项更新放在正确位置。设备章至少要记住一句话：DMA completion 说明设备不再需要某个 buffer，持久化 completion 才说明恢复路径能相信某个状态。

\section{设备模型、IOMMU 与虚拟设备}

\subsection{设备模型：bus、driver、device、class}

Linux 设备模型把发现、匹配、生命周期和 sysfs 统一起来。总线负责枚举和匹配，driver 提供 probe/remove，device 表示具体硬件或虚拟设备，class 给用户空间暴露类别视图，kobject/kset 负责引用计数和 sysfs 层次。

\begin{lstlisting}
struct bus_type      -> match(device, driver)
struct device_driver -> probe(device), remove(device)
struct device        -> parent, bus, driver, kobject
struct class         -> /sys/class view
\end{lstlisting}

probe 成功后驱动必须拥有明确资源集合：MMIO 映射、IRQ、DMA mask、队列、字符/块/网络注册对象。remove 路径按相反顺序撤销，并拒绝新请求、等待 in-flight 请求完成。

\subsection{PCI 驱动骨架}

\begin{lstlisting}
static const struct pci_device_id ids[] = {
  { PCI_DEVICE(VENDOR, DEVICE) },
  { 0 }
};

probe(pdev, id):
  pci_enable_device(pdev)
  pci_request_regions(pdev)
  mmio = pci_iomap(pdev, bar)
  dma_set_mask_and_coherent(pdev, DMA_BIT_MASK(64))
  request_irq(pdev->irq, handler, IRQF_SHARED, name, dev)
  init_queues(dev)

remove(pdev):
  stop_queues(dev)
  free_irq(pdev->irq, dev)
  pci_iounmap(pdev, mmio)
  pci_release_regions(pdev)
  pci_disable_device(pdev)
\end{lstlisting}

错误路径要能从任一步失败回滚。例如 \code{request\_irq} 失败时，要释放 MMIO 和 PCI region；队列初始化失败时，要释放 IRQ。真实源码中这类 \code{goto err\_*} 标签比正常路径更能体现驱动质量。

\subsection{字符、块、网络三类设备}

\begin{longtable}{p{0.18\linewidth}p{0.34\linewidth}p{0.34\linewidth}}
\toprule
类别 & 核心对象 & 典型回调 \\
\midrule
字符设备 & \code{cdev}、\code{file\_operations} & \code{open/read/write/ioctl/poll} \\
块设备 & \code{gendisk}、\code{request\_queue}、blk-mq & \code{submit\_bio}、queue\_rq、completion \\
网络设备 & \code{net\_device}、\code{net\_device\_ops}、NAPI & \code{ndo\_start\_xmit}、poll、set\_rx\_mode \\
\bottomrule
\end{longtable}

同一个硬件可能暴露多个接口。比如 NVMe 是块设备，网卡是网络设备，串口是字符设备。VFS 统一 fd 入口，但驱动对象和 completion 路径完全不同。

\subsection{streaming DMA 与 coherent DMA}

\code{dma\_alloc\_coherent} 返回 CPU 和设备一致可见的内存，适合描述符环；\code{dma\_map\_single/sg} 把已有 buffer 临时映射给设备，适合数据包或块 I/O。streaming DMA 要在方向、同步和 unmap 上严格配对。

\begin{lstlisting}
tx_submit(buf):
  dma = dma_map_single(dev, buf.data, buf.len, DMA_TO_DEVICE)
  if dma_mapping_error(dev, dma): return -EIO
  fill_desc(dma, buf.len)
  dma_wmb()
  ring_doorbell()

tx_complete(desc):
  dma_unmap_single(dev, desc.dma, desc.len, DMA_TO_DEVICE)
  free_buf(desc.buf)
\end{lstlisting}

IOMMU 为设备建立地址空间，限制设备只能 DMA 到授权页。没有 IOMMU 或映射错误时，设备 bug 可能覆盖任意内存；有 IOMMU 时，错误 DMA 会变成 IOMMU fault，至少能阻止破坏扩大。

\subsection{NAPI、NVMe 与 virtio 的共同队列模型}

网络设备、NVMe 与 virtio 的具体队列格式不同，但都可以按同一组问题分析：

\begin{enumerate}
\tightlist
\item probe 如何发现设备、分配队列、注册上层对象。
\item submit 路径如何填描述符、DMA map、通知设备。
\item completion 路径如何从中断/NAPI/poll 回收资源。
\item reset/remove 路径如何停止新请求并清理 in-flight。
\item 错误统计在哪里暴露给 ethtool、sysfs、dmesg 或 block stats。
\end{enumerate}

\subsection{ACPI platform 设备与 PCIe 自描述设备}

并非所有设备都能像 PCIe 一样完整自描述。x86-64 PC/服务器上，可枚举高速设备通常来自 PCIe 配置空间；板载控制器、固件虚拟设备、电源和热管理资源常由 ACPI 表描述。platform driver 的匹配不依赖 vendor/device id，而依赖 ACPI id 或平台总线注册信息。

\begin{lstlisting}
ACPI resource sketch:
  HID/UID identify the platform device
  MMIO resource describes register window
  IRQ resource describes interrupt routing
  _CRS reports current resource settings
  _PS0/_PS3 may describe power-state transitions
\end{lstlisting}

probe 阶段要把描述转换成内核资源：\code{platform\_get\_resource} 取得 MMIO 区间，\code{devm\_ioremap\_resource} 建立映射，\code{platform\_get\_irq} 取得中断，电源管理回调处理设备 suspend/resume。教学系统可以把 ACPI 资源简化成静态设备表，但仍要保留“描述硬件”和“驱动绑定”两个阶段。

\subsection{托管资源和错误回滚}

Linux 驱动大量使用 devm 托管资源。资源绑定到 device 生命周期，remove 或 probe 失败时自动释放。托管资源不是为了偷懒，而是为了减少 probe 多阶段失败路径中的泄漏。

\begin{lstlisting}
probe(dev):
  state = devm_kzalloc(dev, sizeof(*state))
  mmio = devm_ioremap_resource(dev, res)
  irq = platform_get_irq(dev, 0)
  ret = devm_request_threaded_irq(dev, irq, top, thread, flags, name, state)
  if ret:
    return ret
  ret = register_upper_object(state)
  if ret:
    return ret
  return 0
\end{lstlisting}

但 devm 不能替代所有生命周期管理。已经暴露给上层的字符设备、网络设备或块设备必须先停止新请求、撤销注册、等待 in-flight 完成，再让托管资源释放。否则用户态仍可能通过旧 fd 进入已经释放的 MMIO 或队列。

\subsection{ioctl、sysfs 与用户态控制面}

驱动给用户态暴露控制面时，常见三种方式：ioctl、sysfs 属性和 netlink。ioctl 适合和 fd 绑定的设备私有命令；sysfs 适合稳定、简单、可文本化的设备属性；netlink 适合网络和复杂事件通知。控制面设计必须有版本、长度、权限和兼容性边界。

\begin{lstlisting}
ioctl(file, cmd, user_arg):
  switch cmd:
    case GET_STATUS:
      status = snapshot_device_status()
      return copy_to_user(user_arg, status)
    case RESET_DEVICE:
      if !capable(CAP_SYS_ADMIN):
        return -EPERM
      return reset_device_safely()
    default:
      return -ENOTTY
\end{lstlisting}

ioctl 的常见漏洞来自信任用户结构体长度、未检查 padding、把内核指针泄露给用户、或在 reset 中没有和 I/O 路径同步。控制命令不是旁路，它和 read/write/interrupt 一样会并发访问设备状态。

\subsection{runtime PM、suspend 与 resume}

现代设备需要省电。runtime power management 允许设备空闲时关闭时钟或电源域；系统 suspend/resume 要在整机睡眠和唤醒时保存恢复设备状态。驱动必须区分“设备逻辑对象仍存在”和“硬件暂时不可访问”。

\begin{lstlisting}
runtime_suspend(dev):
  stop_new_io(dev)
  wait_inflight(dev)
  save_registers(dev)
  disable_irq(dev)
  disable_clock(dev)

runtime_resume(dev):
  enable_clock(dev)
  restore_registers(dev)
  enable_irq(dev)
  restart_queues(dev)
\end{lstlisting}

如果 resume 后忘记恢复 DMA mask、队列基址或中断 mask，设备可能表现为偶发超时。电源管理 bug 常常只在压力、热插拔、suspend 循环或低功耗平台上出现，因此测试必须覆盖这些状态转换。

\subsection{驱动并发不变量}

驱动至少要维护下面不变量：

\begin{itemize}
\tightlist
\item 提交给设备的描述符在 completion 前不能释放。
\item 设备可 DMA 的 buffer 必须有有效映射，方向必须匹配。
\item remove/reset 开始后，不再接受新请求。
\item 中断处理看到 dying 状态时只确认和丢弃，不再提交新工作。
\item 用户 fd 持有期间，驱动私有对象引用不能降到零。
\item 上半部只做不可睡眠操作，复杂恢复进入线程化 IRQ 或 workqueue。
\end{itemize}

把这些不变量写进测试，才能发现“正常收发能跑”之外的问题。驱动不是能读写设备就合格，而是要在错误、中断风暴、热拔、reset 和并发 close 下仍然收敛。
